\documentclass[12pt]{article}
\usepackage[a4paper,margin=1in]{geometry}\usepackage[T1]{fontenc}
\usepackage{lmodern,microtype}\usepackage{amsmath,amssymb,amsthm,mathtools}
\usepackage{booktabs,array,enumitem,xcolor}\usepackage[numbers,sort&compress]{natbib}
\usepackage[colorlinks=true,linkcolor=blue!55!black,citecolor=blue!55!black,urlcolor=blue!55!black]{hyperref}
\linespread{1.07}
\newtheorem{theorem}{Theorem}[section]\newtheorem{proposition}[theorem]{Proposition}
\newtheorem{lemma}[theorem]{Lemma}
\theoremstyle{remark}\newtheorem{remark}[theorem]{Remark}
\newcommand{\norm}[1]{\left\lVert#1\right\rVert}

\title{An Intrinsic Normalization Audit for the Physical Quartic Divisor\\
\large Eight Predeclared Levels and a Sixteen-Level Frozen Atlas}
\author{Liang Wang}\date{July 2026}
\begin{document}\maketitle
\begin{abstract}
The divisor-first pivot of RH-210--RH-211 leaves a concrete question: can a
simple intrinsic normalization turn the finite physical quartic into a
contracting coefficient flow?  We compare three representations on eight
predeclared small-noise levels and both physical channels: the raw monic
quartic, determinant-radius normalization, and centered root-mean-square
normalization.  No parameters are fitted level by level.

The answer is negative for the two proposed normalizations.  With adjacent
error measured relative to the fine coefficient vector, the raw flow has
range $0.04419$--$0.18945$ and mean $0.10661$.  Determinant-radius
normalization worsens these values to $0.06062$--$0.36787$ and $0.19322$;
centered-RMS normalization gives $0.05941$--$0.23389$ and $0.16592$.
Maximum left/right relative discrepancies are respectively $0.008112$,
$0.019881$, and $0.009368$.  Thus the raw coefficients remain the most
stable of the three tested representations.

We also freeze a sixteen-level dual-channel root atlas down to
$\sigma=0.00125$.  Sparse modulus-cloud extraction exactly reproduces the
previous dense-spectrum anchors while avoiding repeated dense eigensolves.
The negative contraction result is useful: centered-RMS normalization exposes
an exact lower-dimensional shape structure, developed in RH-213, rather than
another fitted scalar law.  No coefficient limit, growing determinant,
Fredholm realization, or Gate-A closure is claimed.
\end{abstract}

\section{Question inherited from the transport frontier}

RH-207 identified a coherent left/right quartic divisor, while RH-202--RH-210
showed that naive state-space transport is not its stable realization
\citep{WangRH207,WangRH210,WangRH211}.  Let
\begin{equation}\label{eq:raw}
 D_{\sigma,s}(z)=\prod_{j=1}^{4}(z-\lambda_{\sigma,s,j})
 =z^4+c_{1,\sigma,s}z^3+c_{2,\sigma,s}z^2
   +c_{3,\sigma,s}z+c_{4,\sigma,s},
\end{equation}
where $s\in\{L,R\}$ labels the two physical channels.  Three coarse anchors
showed channel coherence but not stationarity in $\sigma$.

The present audit asks a deliberately narrow question.  Is the scale motion
mostly an affine gauge that can be removed by one natural formula at every
level?  A positive answer would produce a candidate compact coefficient
family.  A negative answer rules out only the tested formulas, not every
renormalization.

\section{Predeclared scale protocol}

The normalization comparison uses
\begin{equation}\label{eq:auditscales}
 0.04,\ 0.032,\ 0.025,\ 0.02,\ 0.016,\ 0.0125,\ 0.01,\ 0.008.
\end{equation}
The frozen atlas additionally contains
\begin{equation}\label{eq:extrascales}
 0.00625,\ 0.005,\ 0.004,\ 0.0032,\ 0.0025,\ 0.002,\ 0.0016,\ 0.00125.
\end{equation}
The last eight points are not used to redefine the RH-212 winner.  They are
published so that subsequent papers can state their own training and holdout
rules without recomputing the matrix spectrum.

At each scale we retain the same inherited physical rule: remove the Perron
and negative parity modes and select the next four eigenvalues by modulus.
The selected set is checked for conjugate closure.  All 32 endpoint packets
are conjugate closed to the reported floating precision.

\section{Three intrinsic representations}

\subsection{Raw monic coefficients}

The baseline is simply the coefficient vector
\begin{equation}
 C^{\rm raw}_{\sigma,s}=(1,c_1,c_2,c_3,c_4).
\end{equation}
It retains location, radius, and shape.  It is similarity invariant but not
invariant under affine changes of the spectral variable.

\subsection{Determinant-radius normalization}

Set
\begin{equation}\label{eq:detradius}
 \rho=\left|\prod_{j=1}^{4}\lambda_j\right|^{1/4}
      =|c_4|^{1/4},\qquad \widehat\lambda_j=\lambda_j/\rho.
\end{equation}
Then the constant coefficient of
$\prod_j(z-\widehat\lambda_j)$ has modulus one.

\begin{proposition}\label{prop:detintrinsic}
For every quartet with nonzero product, $\rho$ is invariant under root
permutation and operator similarity.  The normalized constant term has
modulus one.
\end{proposition}
\begin{proof}
Both statements follow from the elementary symmetric product and
$\prod_j\widehat\lambda_j=\rho^{-4}\prod_j\lambda_j$.
\end{proof}

This normalization does not remove translation.  It can also amplify the
remaining coefficients when $|c_4|$ is small.

\subsection{Centered RMS normalization}

Define
\begin{equation}\label{eq:rms}
 \mu=\frac14\sum_{j=1}^{4}\lambda_j,
 \qquad
 r=\left(\frac14\sum_{j=1}^{4}|\lambda_j-\mu|^2\right)^{1/2},
 \qquad q_j=\frac{\lambda_j-\mu}{r}.
\end{equation}

\begin{proposition}\label{prop:rmsgauge}
For every nonconstant quartet,
\begin{equation}
 \frac14\sum_j q_j=0,
 \qquad \frac14\sum_j|q_j|^2=1.
\end{equation}
The multiset $\{q_j\}$ is invariant under translation
$\lambda_j\mapsto\lambda_j+a$ and positive dilation
$\lambda_j\mapsto b\lambda_j$ with $b>0$.
\end{proposition}
\begin{proof}
Centering gives the first identity and the definition of $r$ gives the
second.  Under translation, the barycenter translates by the same amount.
Under positive dilation, both centered roots and $r$ are multiplied by $b$.
\end{proof}

For conjugate-closed quartets, $\mu$ is real and the normalized polynomial
has real coefficients and zero cubic coefficient.  This exact structural
gain will matter even though contraction fails.

\section{Error functional}

For adjacent coarse/fine scales $\sigma_c>\sigma_f$, write $C_c,C_f$ for
one of the three monic coefficient vectors and define
\begin{equation}\label{eq:error}
 E_f(C_c,C_f)=\frac{\norm{C_f-C_c}_2}{\norm{C_f}_2}.
\end{equation}
The fine denominator is fixed in advance and treats the refined endpoint as
the target.  Maximum absolute coefficient error and the reverse denominator
are also archived, but no conclusion depends on changing the metric after
inspection.

At a fixed scale, channel discrepancy is
\begin{equation}\label{eq:channelerror}
 E_{LR}=\frac{\norm{C_L-C_R}_2}{\norm{C_L}_2}.
\end{equation}

\section{Finite normalization verdict}

The 14 adjacent/channel cases for each representation give:
\begin{center}
\begin{tabular}{lrrrr}
\toprule
representation & min $E_f$ & max $E_f$ & mean $E_f$ & max $E_{LR}$\\
\midrule
raw & $0.04419$ & $0.18945$ & $0.10661$ & $0.008112$\\
determinant radius & $0.06062$ & $0.36787$ & $0.19322$ & $0.019881$\\
centered RMS & $0.05941$ & $0.23389$ & $0.16592$ & $0.009368$\\
\bottomrule
\end{tabular}
\end{center}

\begin{proposition}[Finite negative result]\label{prop:negative}
On the eight predeclared levels and both channels, neither determinant-radius
nor centered-RMS normalization improves the mean, maximum, or minimum
adjacent fine-relative coefficient error simultaneously.  The raw vector has
the smallest mean and maximum, and it also has the smallest worst channel
discrepancy.
\end{proposition}
\begin{proof}
This is the direct comparison in the displayed table, generated from the
frozen endpoint roots by the exact formulas above.
\end{proof}

The word ``negative'' is important.  The data do not establish that every
intrinsic normalization fails.  They reject two concrete formulas as a
uniformly better cross-scale coordinate on this audit.

\section{Why normalization can worsen the flow}

The determinant radius forces one coefficient to a fixed modulus, but the
other coefficients transform with different powers of $\rho^{-1}$.  Small
relative motion in $c_4$ can therefore induce larger motion in lower-order
terms.  The operation removes one scalar without separating translation from
shape.

Centered-RMS normalization removes two affine degrees of freedom.  If the
remaining shape itself moves substantially, the normalized vector must show
that motion more clearly.  Consequently a larger coefficient error is not a
defect of the quotient.  It can be evidence that location and radius were
masking a genuine shape trajectory.

Indeed, every centered-RMS polynomial in the atlas has zero cubic term and
lies, to roundoff, on one algebraic surface in the remaining three
coefficients.  RH-213 proves that this surface is exactly two-dimensional.
Thus RH-212 changes the next question from ``which scalar should be fitted?''
to ``what exact geometry remains after the affine gauge is removed?''

\section{Sparse extraction and reproducibility}

The historical construction formed a dense deflated bulk operator.  Its
outer nontrivial eigenvalues are unchanged if one instead computes a modulus
cloud of the original sparse folded-Gaussian matrix, removes the Perron and
negative real parity roots, and takes the next quartet.  At the old anchors
$0.04,0.02,0.01$, the resulting roots and coefficients reproduce the frozen
dense results to floating roundoff.

This shortcut is not a new spectral theorem.  It is an algebraic consequence
of replacing two selected eigenvalues by zero in the inherited finite model,
combined with a cloud large enough to resolve the radial gap.  The archive
retains dimensions, roots, coefficients, centers, radii, and conjugacy
diagnostics for every endpoint.

\section{Claim boundary and route consequence}

The following remain open:
\begin{enumerate}[leftmargin=2em]
\item existence of either a raw or normalized coefficient limit;
\item stability under all smaller noise levels;
\item a rank-growing cloud with controlled omitted factors;
\item local uniform convergence to a nonzero analytic determinant;
\item a dynamical or Fredholm realization.
\end{enumerate}

Gate A is therefore open.  Gates B--E are untouched.  No self-adjoint
operator, $T\log T$ law, arithmetic trace identity, zeta divisor, or RH
statement is inferred.  The immediate route is the exact centered shape
manifold, not another post-hoc normalization.

\bibliographystyle{plainnat}\bibliography{references}\end{document}
